\section{The Apollo Guidance Computer}

The Apollo Guidance Computer (AGC) was an early digital computer developed in the 1960s
to help navigate and control Apollo spacecraft on their missions to the Moon. It
provided real-time guidance and control, making it essential for the success of the
Apollo program.

The AGC ran at a speed of 2.048 MHz, which translated to about 40,000 instructions per
second. The AGC had a 16-bit \texttt{word} length, 15 bits for data, and one parity
bit. It had 2 kilowords or 4KB of erasable magnetic-core memory (RAM) along with 36
kilowords or 72KB of fixed core rope memory (ROM) for storing software and mission
data. Astronauts interacted with the AGC using the Display and Keyboard (DSKY), a
21-digit display with a 19-button keyboard that allowed them to input commands and
receive information. This interface was critical for making manual adjustments and
monitoring the system during missions.

Due to these constraints, the operating system, called the Executive, had to
efficiently manage resources. It ensured that important tasks were prioritized over
less critical ones and that key functions could continue even if the system experienced
a reset due to an error.

The Executive was written in AGC assembly, where each instruction was 16 bits long.
Because the instruction set was small, the system also included a sophisticated
interpreter that acted as a virtual machine, providing more advanced
pseudo-instructions beyond the basic AGC instructions.

The details of the AGC Executive explained below is taken from \citet{agc_book}.

\subsection{The Executive} 

The central part of the AGC software was the Executive, that ensured that processes and
tasks got a fair share of the system resources like cpu time and dynamic memory. To
seemlessly perform the real-time control and navigation tasks with extreme reliability
and efficient, the Executive implemented a \textbf{priority-based cooperative
scheduling} algorithm. The algorithm is described below:

\begin{emptythm}
    \begin{itemize}[left=8pt]
        \item The highest running process capitalizes the CPU
        \item Lower priority processes will wait until the highest priority process
            either
            \begin{itemize}
                \item lowers its priority, or
                \item goes to \texttt{sleep} while waiting for some event to happen
            \end{itemize}
        \item While the current highest running process is running, if a new process
            with a higher priority gets scheduled,
            \begin{itemize}
                \item the new process waits until the previous process lowers its
                    priority or goes to \texttt{sleep},
                \item then the new process starts running
            \end{itemize}
    \end{itemize}
\end{emptythm}

\paragraph*{Core sets} Information about the running processes were stored in a table
called Core Sets, which is basically what we call process descriptors
nowadays. Each core set had a structure similar to (but not exactly):

\begin{lstlisting}[style=mycpp, morekeywords={word}]
struct CORE_SET {
    word priority;
    word program_counter; 
    word bbank_reg; // area in unswitched memory where the code is
    word vac_ptr;   // pointer to extra memory

    word mpac[7];   // multi-purpose accumulators
    word mpac_mode; 
};
\end{lstlisting}

Each running process must occupy one of the core sets. There were
\begin{itemize}
    \item 6 available core sets in the Command module, and
    \item 7 available core sets in the Lunar module.
\end{itemize}
In either table, the core set 0 contained the currently running process. So a context
switch basically was equivalent to swapping out a core set with the core set 0.

\paragraph*{Addressable memory} The 7 \texttt{mpac} words had two purposes. In the
normal mode, they gave 7 words of free memory to the programmer to use as they saw fit.
In the interpreter mode, they worked just like accumulator registers, used to store
values of different calculations by the interpreter. If a process needed more memory,
it could request for 42 extra words from the \textbf{Vector Accumulator} area. 

\paragraph*{Starting new process} To start a new process, an empty core set must be
found first, followed by finding a VAC area if requested. After succesfully populating
the core set, the Executive must decide if the new process has a higher priority. 

A special register \texttt{NEWJOB} is used just for this purpose. 
\begin{itemize}
    \item If the value of \texttt{NEWJOB} is 0, then the currently running process has
        the highest priority,
    \item if the value is non-zero, say \texttt{d}, then the process represented by the
        core set at \texttt{d} is the highest priority process.
\end{itemize}
So when a new process is created, its priority is compared with the priority of the
current process, and the \texttt{NEWJOB} register is updated accordingly.

\paragraph*{Switching processes} A process context switch only occurs when the current
process gives up execution, given that the current process continues to probe
\texttt{NEWJOB} every 20ms for a higher priority process. If the \texttt{NEWJOB}
register is not checked for over 640ms, it is assumed that the current process has hung
up, and a hardware restart must be issued to deal with the hangup.

To switch the process, the Executive just needs to store the current context in the
core set at 0, and swap it with the core set at the value stored in \texttt{NEWJOB}.
Since the program counter is read from the core set directly, this automatically
resumes the switched in process from it's last instruction.

If the current process is still the highest running process, and it needs to go to
sleep, or it decreases it's priority, a scan of the core sets table is made to find the
next highest priority process, followed by a switch.

\paragraph*{Sleep and Wakeup} Processes can voluntarily give up execution as it waits
for an external event to occur by going to sleep. The \texttt{JOBSLEEP} instruction
puts the currently running process to sleep. An asleep process is represented by a
negative priority. So all \texttt{JOBSLEEP} needs to do is negate the currently
running processes' priority, and find the next highest priority job.

On the other hand, once the awaited even occurs, another process can wake up a sleeping
process using the \texttt{JOBWAKE} instruction. This instruction negates the negative
priority of the asleep process, which marks it as runnable. It then checks with the
\texttt{NEWJOB} register to decide the highest priority process.

A process can call \texttt{DELAYJOB} instruction to go to sleep for a certain amount of
time. To achieve this, the instruction schedules an ``alarm'' task using the
\textbf{waitlist}, which is described below.


\subsubsection{The Waitlist}

The core processes in the Executive are all co-operative, and even the highest priority
process won't preempt the currently running process until it gives up the execution.
Waitlist provides preemptive scheduling capabilities to the AGC. Tasks that need to be
done within a particular deadline can be scheduled using the waitlist. Every 10ms, a
timer interrupt is generated, which runs any waitlist task that needs to be run.

In essence, the waitlist is a \emph{sorted linked list} of tasks in increasing order of
their deadline. When a new waitlist task \(t\) is scheduled with a particular deadline,
it is inserted into the linked list at the correct position so that
\begin{itemize}
    \item any task \(t_0\) with a deadline less than \(t\) appears before \(t\) in the
        list, and
    \item any task \(t_1\) with a deadline greater than \(t\) appears after.
\end{itemize}
Every 10ms, the timer interrupt checks if the deadline of the first task in the list is
about to expire, and runs that task if so, and then removes the task from the list.

The implementation of waitlist in the Apollo-RTOS is modeled completely after the
implementation in the AGC, so the implemenation details will be discussed later.


\subsection{Interpreter}

\begin{figure}[H]
    \centering
    \includegraphics[width=.5\linewidth]{dsky}
    \caption{DSKY interface diagram by Oona Räisänen}
\end{figure}


\subsection{Design philosophy of AGC}

Unlike most modern general-purpose computers, the AGC operated on a tightly defined
world of hardware and software, where every part of the system could be ``trusted'', in
the sense that every part would behave as expected. As a result, it was deemed
unnecessary to truly isolate the hardware and the OS from the application code. Because
of this, the operating system could make many assumptions limiting the scope of
services, limiting the complexity in design. This resulted in a small but efficient
operating system. 

This theme of mutual trust within the operating system will be prevelant throughout the
project. Since the scope of a microcontroller is quite minimal, and resources are
limited, it makes sense to not make any kernel/user space isolation. All usages of
system resources is already anticipated during development. However, to minimize
possibility of programming errors, it's necessary to enclose system resources under
system and library calls to minimize human programming errors.

\begin{figure}[H]
    \centering \includegraphics[width=.6\linewidth]{margaret_hamilton}
    \caption{Margaret Hamilton standing next to listings of the software she and her
    MIT team produced for the Apollo Project.}
\end{figure}

\subsubsection{Design consequences}
There are a few consequences of this design of the scheduler:

Since all the processes use the same address space, and inter process communication is
done through shared memory

Letting the currently running process decide when to give up the execution context
fixes several nasty race conditions usually prevelant in preemptive schedulers. For
example,
\begin{itemize}
    \item In multi-threaded pro
\end{itemize}

